% 分野別類題: 1階線形非同次微分方程式 解答例
% コンパイル: TEXINPUTS="../../../00_common:$TEXINPUTS" latexmk -xelatex answers.tex
\documentclass[a4paper,11pt]{article}
\usepackage{preamble}
\usepackage{shoakubox}

\begin{document}

\kamokumei{類題演習 解答例 --- 1階線形非同次微分方程式（積分因子法）}

\daimon{【解法】1階線形非同次微分方程式の解き方と導出}

\noindent
1階線形非同次微分方程式とは、次の形の方程式である。
\[
\frac{dy}{dx} + P(x)\,y = Q(x)
\]
これを解くために、積分因子（積分因数）
\[
\mu(x) = \exp\left(\int P(x)\,dx\right)
\]
を両辺に掛ける。このとき
\[
\mu(x)\frac{dy}{dx} + \mu(x) P(x)\,y = \mu(x)\,Q(x)
\]
であり、$\mu'(x) = \mu(x) P(x)$ に注意すると、左辺は積の微分法則により
\[
\mu(x)\frac{dy}{dx} + \mu'(x)\,y = \frac{d}{dx}\bigl(\mu(x)\,y\bigr)
\]
とまとめられる。したがって元の方程式は
\[
\frac{d}{dx}\bigl(\mu(x)\,y\bigr) = \mu(x)\,Q(x)
\]
という「両辺が完全微分の形」に帰着する。両辺を $x$ で積分すれば
\[
\mu(x)\,y = \int \mu(x)\,Q(x)\,dx + C
\]
となり、これより
\[
y = \frac{1}{\mu(x)}\left(\int \mu(x)\,Q(x)\,dx + C\right)
\]
が一般解として得られる（$C$ は任意定数）。初期条件が与えられた場合は、一般解に代入して $C$ を定める。なお $Q(x) = 0$ の場合はこの解法は前分野の一次同次線形微分方程式（変数分離）に一致し、$y = C/\mu(x) = C\exp\left(-\int P\,dx\right)$ となる。

\clearpage
\begin{mondaiboxS}{問1}
$\dfrac{dy}{dx} + \dfrac{y}{x} = \dfrac{\sin x}{x} \quad (x > 0)$ の一般解を求めよ。
\end{mondaiboxS}
\kaitou
$P(x) = 1/x$ であるから、積分因子は
\[
\mu(x) = \exp\left(\int \frac{1}{x}\,dx\right) = \exp(\ln x) = x
\]
両辺に $\mu(x) = x$ を掛けると
\[
x\frac{dy}{dx} + y = \sin x
\quad\Longleftrightarrow\quad
\frac{d}{dx}(xy) = \sin x
\]
両辺を積分して
\[
xy = -\cos x + C
\]
よって
\[
\boxed{\; y = \frac{C - \cos x}{x} \;}
\qquad (C \text{ は任意定数})
\]
（検算: $y = (C-\cos x)/x$ を微分すると $y' = \dfrac{\sin x \cdot x - (C-\cos x)}{x^{2}} = \dfrac{\sin x}{x} - \dfrac{C-\cos x}{x^{2}} = \dfrac{\sin x}{x} - \dfrac{y}{x}$。すなわち $y' + y/x = \sin x / x$ が成り立つ。）

\clearpage
\begin{mondaiboxS}{問2}
$\dfrac{dy}{dx} - 2x\,y = e^{x^{2}}$ の一般解を求めよ。
\end{mondaiboxS}
\kaitou
$P(x) = -2x$ であるから、積分因子は
\[
\mu(x) = \exp\left(\int (-2x)\,dx\right) = e^{-x^{2}}
\]
両辺に $\mu(x) = e^{-x^{2}}$ を掛けると
\[
e^{-x^{2}}\frac{dy}{dx} - 2x\,e^{-x^{2}}y = e^{-x^{2}}e^{x^{2}} = 1
\quad\Longleftrightarrow\quad
\frac{d}{dx}\bigl(e^{-x^{2}}y\bigr) = 1
\]
両辺を積分して
\[
e^{-x^{2}}y = x + C
\]
よって
\[
\boxed{\; y = (x + C)\,e^{x^{2}} \;}
\qquad (C \text{ は任意定数})
\]
（検算: $y = (x+C)e^{x^{2}}$ を微分すると $y' = e^{x^{2}} + (x+C)\cdot 2x\,e^{x^{2}} = e^{x^{2}}\bigl(1 + 2x^{2} + 2xC\bigr)$。一方 $2xy = 2x(x+C)e^{x^{2}} = e^{x^{2}}(2x^{2}+2xC)$ であるから、$y' - 2xy = e^{x^{2}}\bigl(1+2x^{2}+2xC-2x^{2}-2xC\bigr) = e^{x^{2}}$ が成り立つ。）

\clearpage
\begin{mondaiboxS}{問3}
$\dfrac{dy}{dx} + 2y = e^{-x}, \quad y(0) = 3$ を解け。
\end{mondaiboxS}
\kaitou
$P(x) = 2$ であるから、積分因子は
\[
\mu(x) = \exp\left(\int 2\,dx\right) = e^{2x}
\]
両辺に $\mu(x) = e^{2x}$ を掛けると
\[
e^{2x}\frac{dy}{dx} + 2e^{2x}y = e^{2x}e^{-x} = e^{x}
\quad\Longleftrightarrow\quad
\frac{d}{dx}\bigl(e^{2x}y\bigr) = e^{x}
\]
両辺を積分して
\[
e^{2x}y = e^{x} + C
\quad\Longrightarrow\quad
y = e^{-x} + Ce^{-2x}
\]
初期条件 $y(0) = 3$ より $1 + C = 3$、すなわち $C = 2$。したがって
\[
\boxed{\; y = e^{-x} + 2e^{-2x} \;}
\]
（検算: $y' = -e^{-x} - 4e^{-2x}$ であり、$2y = 2e^{-x} + 4e^{-2x}$ であるから、$y' + 2y = (-e^{-x}-4e^{-2x}) + (2e^{-x}+4e^{-2x}) = e^{-x}$ が成り立つ。また $y(0) = 1 + 2 = 3$ も満たす。）

\end{document}
